\section{Limitations}

CLOC is a formulation of preservation, not a claim that all old-task behavior can be retained. Per-sample equality is stronger than expected-risk preservation and can conflict with fitting later tasks. The penalty coefficient therefore controls a stability--plasticity relaxation rather than guaranteeing feasibility of the original hard constraint.

IEWC, like EWC, compresses an old task into one importance matrix evaluated at $\theta_A^\star$ and keeps it fixed during later training. Before aggregation, however, its normalized parameter-sensitivity direction $a^{G}(x;\theta_A^\star)=J_{\theta_A^\star}(x)^\top\hat r^{G}(x)$ is input-dependent. Circuit analyses make the same distinction between a model's fixed connections and the input-dependent feature contributions that realize a particular computation \citep{DunefskyTranscoders}. A single average---especially after diagonalization---can therefore hide sensitivity confined to subsets of inputs or distributed across overlapping pathways, and cannot track changes in the relevant Jacobians as the parameters move away from $\theta_A^\star$. Sample weighting changes which empirical average is stored, but cannot recover this conditional structure after compression. Conditional mixtures or dynamically refreshed importance are natural extensions, at the cost of richer stored summaries or access to old inputs.

The output metric must be specified as a matched primal/dual pair. Ambient singular metrics require restriction to an identifiable subspace and a stated pseudoinverse convention. Our ordered-CDF construction satisfies that requirement on the probability-simplex tangent, but the inherited diffusion result is exploratory and does not validate a general Wasserstein geometry.

Practical IEWC also introduces approximations beyond the derivation. A diagonal retains coordinatewise importance but discards correlations, eigenspaces, and reparameterization behavior. Importance estimates use finitely many old samples. RDW is an order-dependent weighting heuristic with additional cost and is not implied by level-set normalization. Relative damping guarantees invariance to a global change of loss units, not to arbitrary sample-dependent transformations at finite damping.

Finally, empirical comparisons cover selected datasets, model families, adaptation regimes, and task orders. Validation selection and locked paired evaluation reduce selection bias but cannot establish a domain-independent advantage or deployment robustness. Offline/joint training is only a reference because it violates the continual data-access constraint, and untested replay, distillation, and parameter-isolation methods may be preferable when their storage or architectural assumptions are allowed.
